\chapter{Mar.~3 --- Differential Forms}

\section{Differential Forms}

\begin{definition}
  For a variety $X$ and
  $0 \le p \le \dim X$,
  \[
    \Omega_X^p := \wedge^p \Omega_X
  \]
  is the space of \emph{$p$-forms}
  on $X$. If $X$ is smooth and
  irreducible, its \emph{canonical bundle}
  is
  \[
    \omega_X := \wedge^{\dim X} \Omega_X
    = \det(\Omega_X).
  \]
\end{definition}

\begin{remark}
  If $M$ is a smooth compact oriented
  manifold, the \emph{de Rham cohomology}
  is given by
  \[
    H_{\mathrm{dR}}^k(M, \R)
    = \frac{\text{closed $k$-forms}}{\text{exact $k$-forms}}.
  \]
  Then Poincar\'e duality says that
  (here $n = \dim X$)
  \begin{align*}
    H_{\mathrm{dR}}^k(M, \R) \times H_{\mathrm{dR}}^{n - k}(M, \R)
    &\longrightarrow H_{\mathrm{dR}}^n(M, \R) \cong \R \\
    ([\alpha], [\beta]) &\longmapsto [\alpha \wedge \beta]
  \end{align*}
  is a perfect pairing, i.e. the following
  map is an isomorphism:
  \begin{align*}
    H_{\mathrm{dR}}^k(M, \R) &\longrightarrow H_{\mathrm{dR}}^{n - k}(M, \R)^\vee \\
    [\alpha] &\longmapsto ([\beta] \mapsto [\alpha \wedge \beta])
  \end{align*}
  The canonical bundle $\omega_X$ will be an
  analogue of
  $H_{\mathrm{dR}}^n(M, \R) \cong \R$ in
  algebraic geometry.
\end{remark}

\section{Canonical Divisor}
\begin{definition}
  The \emph{canonical divisor} of a smooth
  irreducible variety $X$ is a divisor
  $K_X$ such that $\omega_X \cong \OO_X(K_X)$ (note that technically
  $K_X$ is only defined up to linear
  equivalence).
\end{definition}

\begin{example}
  How do we compute the canonical divisor
  $K_X$?
  \begin{enumerate}
    \item First consider $\Affine^n$.
      In this case $\Omega_{\Affine^n} = \OO_{\Affine^n}^{\oplus n}$,
      so $\det(\Omega_{\Affine^n}) = \OO_{\Affine^n}$.
      So $K_{\Affine^n} = 0$ (alternatively
      one can note that $\Cl(\Affine^n) = 0$,
      so there every divisor is $0$ up
      to linear equivalence).
    \item Next consider $\PP^n$. We
      need to use the following tools:
      \begin{enumerate}[i)]
        \item The Euler short exact sequence
          $0 \longrightarrow \Omega_{\PP^n} \longrightarrow \OO_{\PP^n}(-1)^{\oplus (n + 1)} \longrightarrow \OO_{\PP^n} \longrightarrow 0$.
        \item If $0 \longrightarrow \mathcal{E} \longrightarrow \mathcal{F} \longrightarrow \mathcal{G} \longrightarrow 0$ is a short exact sequence of locally free sheaves
          of ranks $e, f, g$ on $X$, then
          $\det(\mathcal{F}) \cong \det(\mathcal{E}) \otimes \det(\mathcal{G})$.

          \pagebreak
          To see this, note that given
          a short exact sequence of
          free $R$-modules
          \[
            \begin{tikzcd}
              0 \ar[r] & E \ar[r] & F \ar[r] & G \ar[r] & 0
            \end{tikzcd}
          \]
          of ranks $e, f, g$ for some
          ring $R$,
          we can define a map
          \begin{align*}
            \det(E) \otimes \det(G) &\longrightarrow \det(F) \\
            (v_1 \wedge \cdots \wedge v_e) \otimes (\overline{w}_1 \wedge \cdots \wedge \overline{w}_g) &\longmapsto v_1 \wedge \cdots \wedge v_e \wedge w_1 \wedge \cdots \wedge w_g,
          \end{align*}
          where we think of $G = F / E$.
          One can check that this is
          well-defined, is an isomorphism
          (it is enough to check this on
          a basis), and commutes with
          localization. This induces a
          morphism
          \[
            \det(\mathcal{F})|_U
            \longrightarrow
            \det(\mathcal{E})|_U \otimes \det(\mathcal{G})|_U
          \]
          on $U \subseteq X$ open affine
          which trivializes
          $\mathcal{E}, \mathcal{F}, \mathcal{G}$.
          Since the map commutes
          with localization, these morphisms
          glue to give an isomorphism
          $\det(\mathcal{F}) \to \det(\mathcal{E}) \otimes \det(\mathcal{G})$.
        \item For $\mathcal{L}_1, \dots, \mathcal{L}_r$
          invertible sheaves, we have
          \[
            \det(\mathcal{L}_1 \oplus \cdots \oplus \mathcal{L}_r)
            \cong \mathcal{L}_1 \otimes \cdots \otimes \mathcal{L}_r.
          \]
          One can prove this as above
          by defining a map on affine opens
          and showing that it commutes with
          localization, or by checking the
          transition data.
      \end{enumerate}
      Using these three properties, we see
      that
      \[
        \omega_{\PP^n}
        = \det \Omega_{\PP^n}
        = \det(\OO_{\PP^n}(-1)^{\oplus (n + 1)})
        \otimes \det(\OO_{\PP^n})
        = \OO_{\PP^n}(-n - 1).
      \]
      Hence we see that $K_{\PP^n} = -(n + 1) H$
      with $H \subseteq \PP^n$ a
      hyperplane.
  \end{enumerate}
\end{example}

\section{Exact Sequences of Differentials}
\begin{prop}\label{prop:exact-sequences-differentials}
  We have the following:
  \begin{enumerate}
    \item For a morphism of varieties
      $f : X \to Y$, there is an exact sequence
      \[
        \begin{tikzcd}
          f^* \Omega_Y
          \ar[r] & \Omega_X
          \ar[r] & \Omega_{X / Y}
          \ar[r] & 0
        \end{tikzcd}
      \]
    \item For $i : Z \hookrightarrow X$
      a closed embedding, there is an
      exact sequence
      (where $\mathcal{I} = \mathcal{I}_Z \subseteq \OO_X$)
      \[
        \begin{tikzcd}
          \mathcal{I} / \mathcal{I}^2
          \ar[r] & i^* \Omega_X
          \ar[r] & \Omega_Z
          \ar[r] & 0
        \end{tikzcd}
      \]
  \end{enumerate}
\end{prop}

\begin{proof}
  (1) For $U \subseteq X$ open affine with
  $f(U) \subseteq V$ open affine, set
  \[
    A = k \longhookrightarrow B = \OO_Y(V)
    \longhookrightarrow C = \OO_X(U).
  \]
  From this, we get an exact sequence
  \[
    \begin{tikzcd}
      \Omega_{B / A} \otimes_B C
      \ar[r] & \Omega_{C / A}
      \ar[r] & \Omega_{C / B}
      \ar[r] & 0
    \end{tikzcd}
  \]
  Then use that this sequence is natural
  (i.e. for another row $A' \hookrightarrow B' \hookrightarrow C'$ with the
  appropriate arrows, there is another
  row in the exact sequence which
  commutes)
  and commutes with localization
  to get an exact sequence of coherent
  sheaves.

  (2) Reduce to the corresponding
  ring theoretic statement.
\end{proof}

\begin{remark}
  In Proposition \ref{prop:exact-sequences-differentials}(2),
  we abusively write
  $\mathcal{I} / \mathcal{I}^2$ for
  $i_*(\mathcal{I} / \mathcal{I}^2)$, since
  for $X$ affine,
  $A = \OO_X(X)$, and $I = I(Z) \le A$,
  we have that $I / I^2$ is naturally
  an $A / I$-module as $\mathrm{Ann}(I / I^2) \subseteq I$.
  Thus
  \[
    i^*(\mathcal{I} / \mathcal{I}^2)
    = (I / I^2 \otimes_A A / I)^\sim
    = \widetilde{I / I^2}.
  \]
  More generally, there is an equivalence
  of categories
  \[
    \begin{tikzcd}
      \Coh_Z \ar[r, shift left=1, "i_*"] & \text{coherent sheaves $\mathcal{F}$ on $X$ with $\mathrm{Ann}(\mathcal{F}) \subseteq \mathcal{I}$} \ar[l, shift left=1, "i^*"]
    \end{tikzcd}
  \]
\end{remark}

\begin{example}
  For $i : \{p\} \hookrightarrow \PP^1$,
  we have an exact sequence
  \[
    \begin{tikzcd}
      0 \ar[r] & \mathcal{I}_{\{p\}}
      \ar[r] & \OO_{\PP^1}
      \ar[r] & i_* \OO_{\{p\}}
      \ar[r] & 0
    \end{tikzcd}
  \]
\end{example}

\begin{remark}
  If $Z, X$ are smooth varieties of
  pure dimension, then
  \begin{itemize}
    \item Mustata 6.3.21 implies
      that $\mathcal{I} / \mathcal{I}^2$ is a locally free $\OO_X$-module of rank $\dim X - \dim Z$.
    \item We call
      $N_{Z / X} = \Hom_{\OO_Z}(\mathcal{I} / \mathcal{I}^2, \OO_Z)$ the
      \emph{normal bundle}.
    \item We have a short exact sequence:
      \[
        \begin{tikzcd}
          0 \ar[r] & \mathcal{I} / \mathcal{I}^2
          \ar[r] & i^* \Omega_X
          \ar[r] & \Omega_Z
          \ar[r] & 0
        \end{tikzcd}
      \]
      To see this, note that locally
      on $U \subseteq X$ open affine
      trivializing the locally free sheaves,
      we get an exact sequence
      \[
        \begin{tikzcd}
          \OO_U^{\dim X - \dim Z}
          \ar[r, "\alpha"] & \OO_U^{\dim X}
          \ar[r, "\beta"] & \OO_U^{\dim Z}
          \ar[r] & 0
        \end{tikzcd}
      \]
      We want to check that we get
      a short exact sequence if we add a
      $0$ on the left, i.e. that $\ker \alpha = 0$.
      For $x \in U$, we can apply
      the right exact functor
      $- \otimes \OO_U / \m_x\ (\cong k)$ to get
      \[
        \begin{tikzcd}
          k^{\dim X - \dim Z}
          \ar[r] & k^{\dim X}
          \ar[r] & k^{\dim Z}
          \ar[r] & 0
        \end{tikzcd}
      \]
      which is short exact by dimension
      reasons. So we see that
      \[
        (\ker \alpha)_x
        \subseteq \m_x^{\dim X - \dim Z}
        \subseteq \OO_U^{\dim X - \dim Z}.
      \]
      As this holds for all $x \in U$,
      we get $\ker \alpha = 0$. So the
      original sequence is short exact.
  \end{itemize}
\end{remark}

\begin{corollary}
  If $i : Z \hookrightarrow X$ is a closed
  embedding of smooth irreducible varieties,
  then
  \[
    \omega_Z
    = \det(\Omega_Z)
    \cong \det(i^* \Omega_X) \otimes \det(N_{Z / X})
    = i^* \omega_X \otimes \det(N_{Z / X}).
  \]
\end{corollary}

\section{Graded Modules}

\begin{remark}
  On an affine variety $X$, there is a
  bijection
  \begin{align*}
    \{\text{finitely generated $A(X)$-modules}\} / {\cong}
    &\longleftrightarrow
    \{\text{coherent sheaves on $X$}\} / {\cong} \\
    M &\longmapsto \widetilde{M} \\
    \mathcal{F}(X)
    &\mathrel{\reflectbox{\ensuremath{\longmapsto}}}
    \mathcal{F}.
  \end{align*}
  Now we want a similar correspondence
  for a closed subvariety
  $X \subseteq \PP^n$ with
  $S(X) = k[x_0, \dots, x_n] / I(X)$ and
  graded $S(X)$-modules, though we will
  only get a bijection up to some further
  equivalence.
\end{remark}

\begin{remark}
  Let $X \subseteq \PP^n$ be a closed
  subvariety and $S = \bigoplus_{d \in \N} S_d$
  its homogeneous coordinate ring.
\end{remark}

\begin{definition}\label{def:graded-module}
  A \emph{graded $S$-module} $M$ is an
  $S$-module with a direct sum decomposition
  \[
    M = \bigoplus_{e \in \Z} M_e
  \]
  as abelian groups, such that
  $S_d \cdot M_e \subseteq M_{d + e}$.
  A \emph{morphism} of graded $S$-modules
  \[
    \phi : M \longrightarrow N
  \]
  is a morphism of $S$-modules which
  respects the grading, i.e.
  $\phi(M_d) \subseteq N_d$ for all
  $d \in \Z$.
\end{definition}

\begin{example}
  We have the following:
  \begin{enumerate}
    \item If $M, N$ are graded $S$-modules,
      then so are $M \oplus N$ and
      $M \otimes_S N$, with
      \[
        (M \oplus N)_d
        = M_d \oplus N_d \quad \text{and} \quad
        (M \otimes_S N)_d = \bigoplus_{e + f = d} (M_e \otimes_{S_0} N_f).
      \]
    \item For $m \in \Z$, let $S(m)$
      be the graded $S$-module such that
      $S(m)_d = S_{m + d}$.
    \item If $\phi : M \to N$ is a
      morphism of graded $S$-modules, then
      $\ker \phi$, $\im \phi$, $\coker \phi$ are
      graded $S$-modules.
  \end{enumerate}
\end{example}

\begin{exercise}
  If $M$ is a finitely generated
  graded $S$-module, then show that
  there exists
  an exact sequence of graded
  $S$-modules of the form
  \[
    \begin{tikzcd}
      \bigoplus_{i = 1}^s S(a_i)
      \ar[r, "\alpha"] & \bigoplus_{j = 1}^r S(b_j)
      \ar[r, "\beta"] & M \ar[r] & 0
    \end{tikzcd}
  \]
  Hint: Choose generators $m_1, \dots, m_r \in M$
  that are homogeneous, set
  $b_j = - \deg m_j$,
  and let $\beta(e_j) = m_j$.
\end{exercise}

\begin{remark}
  Next time, we will define
  $\widetilde{M} \in \QCoh_X$, and we will
  see that $\widetilde{S(m)} \cong \OO_X(m)$.
\end{remark}
